\documentclass{ltjsbook}
\usepackage{docmute} 
\usepackage{../../myStyle} 

\begin{document}
% \chapter{データ生成モデリング}
\chapter{Data Generation Modeling}
\label{x17_Bayesian}

% \section{データ生成モデリング}
\section{Data Generation Modeling}
% 「心理学統計法」と名のつく講義のほとんどは，記述統計や推測統計，とくに帰無仮説検定を扱うことが中心であるのが現状です\footnote{とくに公認心理士に必要な授業として，科目名「心理学統計法」を名乗る必要があり，そこで教える基本的な内容としてこれらが含まれています。なお測定論やそれに関する多変量解析法はそれほど中心的話題ではありません。みんな使うのになあ。}。心理学において帰無仮説検定を行う理由は，心理学実験の効果を検証すること，それも手元の標本についての記述ではなく母集団においてもその効果があると言えるかどうかを検証するためです。手元のデータに基づいて，目に見えない母集団全体のことに推論するというのは，基本的に不良設定問題です。つまり，確実な正解を導き出すにはヒントが少なすぎる，圧倒的に不利な状況で知恵を絞る必要があるのです。こうした不利な状況下ですから，いくつかの仮定をおいて，確率的に推論するのでした。その推論の方法として，\keyterm{モーメント法}{もーめんとほう}{moment method}，\keyterm{最尤法}{さいゆうほう}{Maximum Likelohood estimation}，\keyterm{ベイズ法}{べいずほう}{Bayesian estimation}があるのでした\footnote{これら3つの方法について，十分に思い出せない人は一年時の授業資料を振り返ってみてください。}。
Most lectures titled "Statistical Methods in Psychology" currently focus on dealing with descriptive statistics and inferential statistics, particularly null hypothesis testing\footnote{Particularly, it is necessary to adopt the course name "Statistical Methods in Psychology" as a required course for qualified psychologists, and these topics are included as the basic content to be taught. Note that measurement theory and related multivariate analysis methods are not so much a central topic. I wonder why everyone uses them.}. The reason for performing null hypothesis testing in psychology is to verify the effect of psychological experiments, not just descriptions of the samples at hand, but to verify whether this effect also applies to the population. Based on the data at hand, inferring about the unseen entire population is fundamentally a poorly-defined problem. In other words, there is too little hint to derive a definite correct answer, and you must rack your brains in an overwhelmingly disadvantageous situation. Under such disadvantageous circumstances, some assumptions were made, and inference was made probabilistically. As a method of this inference, there were the \keyterm{moment method}{moment method}{moment method}, the \keyterm{maximum likelihood method}{maximum likelihood method}{Maximum Likelohood estimation}, and the \keyterm{Bayesian method}{Bayesian method}{Bayesian estimation}\footnote{If you can't remember enough about these three methods, please look back at the first-year course materials.}.

% これらの推定方法が\keyterm{実験計画}{じっけんけいかく}{Experimerntal Design}(要因計画ともいう)と組み合わせで用いられ，帰無仮説検定となるのでした。心理学における実験的なアプローチは，実験群と統制群に無作為に割り当てた集団に対し，介入・処置のあとでの群平均を見ることでその効果を検証します。人間を相手にする研究ですから，当然誤差や個人差がデータには含まれますが，無作為割り当てと平均化によってそれらはキャンセルアウトされ，平均の比較をすることで効果を見ることができると考えるのでした。また標本の平均ではなく，それを用いて母平均を推測することで，結果の一般化を考えます。母平均の推定には点推定と区間推定とがあり，確率的表現を用いた区間推定をつかって慎重に結論を導き出すのです。ただし区間推定は判断基準が明確ではありませんから，帰無仮説と対立仮説という2つのモデルを戦わせて決着を見る，というやり方をして一応の決着を見るのでした。こうした「実験計画」+「推測統計学」＝「帰無仮説検定」は，長らく心理学のスタンダードとして君臨しています。
These estimation methods were used in combination with \keyterm{experimental design}{experimerntal design}, also known as factor planning, and it was a null hypothesis test. The experimental approach in psychology verifies its effects by looking at the group average after intervention and treatment on the group randomly assigned to the experimental group and the control group. This is a study that deals with humans, so of course, errors and individual differences are included in the data, but through random assignment and averaging, they are cancelled out, and the effect can be seen by comparing averages. Also, instead of the sample mean, we consider generalization of the results by using it to estimate the population mean. There are point estimates and interval estimates for estimating the population mean, and we carefully draw conclusions using interval estimates with probabilistic expressions. However, because the criteria for interval estimation are not clear, we settle for a provisional conclusion by pitting two models, the null hypothesis and the alternative hypothesis, against each other to see the outcome. This "experimental design" + "inferential statistics" = "null hypothesis test" has reigned as the standard in psychology for a long time.

% こうしたアプローチについて，昨今批判があることについては，今は触れないでおきましょう\footnote{\textcite{amrhein2019scientists}などが指摘するように，誤った使い方をされることが多く心理学研究の意義そのものが疑われるほどになっている現状があります。また\textcite{Toyoda202003}ではさらに辛辣に批判がなされています。}。それよりもこうした問いの立て方や解決策に注目してみたいと思います。心理学は物理学を範として，科学scienceの仲間入りをしたい，というモチベーションが強く根付いている世界であり，また人間というのは嘘をついたり間違えたりするものだ，ということが骨身に染みてわかっていますから\footnote{というか心理学の研究というのは，いかに人間がダメで偏った考え方をする生き物であるかを，滔々と明らかにしていくという側面もあります。}，データを分析する際にも\keytermL{客観性}{きゃっかんせい}を大事にすることが殊更重視されます。客観性の反対は主観性，つまり本人の思い込みや考え方の癖がもっとも邪魔になるのです。そこでデータに対しては真っ白な気持ちで向き合うものだ，という姿勢をとることになります。
Let's not touch on the criticism that these approaches have received recently\footnote{There is a situation where misuse is common, and the significance of psychological research itself is being questioned, as pointed out by \textcite{amrhein2019scientists}. There is also more scathing criticism in \textcite{Toyoda202003}.}. Instead, I would like to focus on the way these questions are posed and the solutions they propose. Psychology is a field deeply rooted in the motivation to join the ranks of science, following the model of physics, and where it is deeply understood that humans lie and make mistakes\footnote{In other words, psychological research also has a side where it reveals how humans are flawed and have biased ways of thinking.}. Therefore, when analyzing data, it is particularly emphasized to value \keytermL{objectivity}{concrete detail}. The opposite of objectivity is subjectivity, which is where the person's own preconceptions and biases become the biggest obstruction. Therefore, the approach to data calls for adopting an attitude of facing it with a blank mind.

% 言い方を変えると，分析をする前はデータについて何も考えてないよ，という態度をとるわけです。もちろんデータは要因計画の結果として得られるものですから，計画を立てる際は主観的な誤謬が微塵も入り込むことないように緻密に練り上げるのですが，出てきた結果は結果，数字の羅列に過ぎないと考えるのです。その結果を統計的に分析するときは，それが記憶実験の数字であろうが臨床実験の結果であろうが気にすることはなく，淡々と帰無仮説検定の俎上に乗せていくことになります。こうしたアプローチができるからこそ，そしてそれを許す統計ソフトウェアがあるからこそ，心理学の研究を積み重ねていくことができたのだという一面があります。こうした研究アプローチは，データ駆動型分析と言えるでしょう。つまりデータができてから，分析が始まるという考え方です。
To put it another way, before doing an analysis, we take the attitude that we're not thinking about the data at all. Of course, because data is obtained as a result of factor planning, when making a plan, we meticulously work out so that no subjective errors will creep in. However, we consider the results that come out to be nothing more than a string of numbers. When statistically analyzing these results, we don't care whether they are numbers from a memory experiment or results from a clinical experiment, we simply and calmly put them on the chopping block of null hypothesis testing. This approach has been possible, and there is statistical software that allows it, which is one reason why we have been able to accumulate research in psychology. This research approach can perhaps be called data-driven analysis. That is to say, the analysis begins after the data is created.

% これに対して，データがどのように生まれてきたのか，そのメカニズムを考え，その仕組みを明らかにしていこうというのが\keytermL{データ生成モデリング}{でーたせいせいもでりんぐ}です。\textcite{Matsuura201610}は統計モデリングを「確率モデルをデータに当てはめて現象の理解と予測を促す営みのこと」と定義しています。すなわちこのアプローチは，まずデータがどのようなメカニズムで生まれてきたのかを，簡潔な数式を使って表現します。そのモデルをデータに当てはめ，このモデルからデータが出てきたと言えるかどうか，他のモデルの方が今のデータをうまく説明するのではないか，といった比較検証をしていくことになります。データ駆動型分析では，こうしたメカニズムが実験計画の中に暗黙理に埋め込まれていたと言えるかもしれません。データ生成モデル駆動型の分析は，メカニズムを明示して検証する，というところが違います。
In response to this, \keytermL{data generation modeling}{でーたせいせいもでりんぐ} is about considering the mechanism by which data was generated and clarifying that system. \textcite{Matsuura201610} defines statistical modeling as "the act of promoting understanding and prediction of phenomena by applying probability models to data". In other words, this approach starts by expressing how the data was generated using simple equations. Then, you apply that model to the data, comparing and verifying whether the data can be said to have come from this model, or if another model better explains the current data. In data-driven analysis, you could say that such a mechanism is implicitly embedded in the experimental design. The difference with data generation model-driven analysis is that it explicitly tests the mechanism.

% データが作られるメカニズムを考えるアプローチの利点は，予測にも向います。メカニズムが正しい，あるいはうまく現状のデータを再現できるのであれば，おそらく今後も同じようにデータが作られていくでしょう。であれば将来の予測ができる，という考え方です。たとえば市場の動向の予測，消費者の傾向の予測ができればそのメリットは想像に難くありません。昨今はデータサイエンスという言葉が流行していますが，そこにはこうした研究アプローチの隆盛があるのです。もちろん，心理学においてもモデリングのアプローチは有用ですし，従来通りの平均値の比較をする上でも，さまざまなメリットがあります。
The advantage of thinking about the mechanism by which data is created is that it also lends itself to prediction. If the mechanism is correct, or if it can replicate the current data well, then it is likely that data will continue to be generated in the same way in the future. The idea is that if we can predict the future based on this, the benefits are easy to imagine. If you can predict market trends or consumer trends, the benefits are not hard to imagine. The word "data science" is popular these days, and behind it lies the rise of this kind of research approach. Of course, the modeling approach is also useful in psychology, and there are various merits to comparing mean values as usual.

% この授業では，データ生成モデリングのアプローチをとって，心理学的にもどういう意義があるのかを考えていきましょう。
In this class, let's consider the significance from a psychological perspective, taking a data generation modeling approach.

% これを考える上で重要なのが\keyterm{確率モデル}{かくりつもでる}{stochastic model}です。言葉の通り，データ生成過程を確率を使って表現します。なぜ確率を？と思うかもしれませんが，手元のデータは誤差や個人差を含んで微小に変わる，偶然の成分が必ず入っているからです。こうした偶然をハンドリングし，偶然の中にも理論的な予測をするという点で確率が必要になるのです。また確率モデルで考えるときに，最尤法よりもベイズ法の方がよく使われます。確率モデルが複雑になっていくにつれ，最尤法では計算コストが非常に高く，実質的に解が得られないということも少なくありません。ベイズ法によるアプローチはこの問題をクリアしてくれるからです。それでは改めて，ベイズ法について復習しておきましょう。
One important consideration is the \keyterm{stochastic model}. As the name implies, it represents the data generation process using probability. You may be wondering why use probability? This is because the data at hand always contains a chance component, such as errors and individual differences, that can vary slightly. Probability is necessary in order to handle these chance occurrences and to make theoretical predictions within these uncertainties. When thinking in terms of a probability model, the Bayesian method is often used more than the maximum likelihood method. As the probability model becomes more complex, the calculation cost in the maximum likelihood method can be very high, and it's not uncommon to be effectively unable to obtain a solution. This is a problem that the Bayesian approach can overcome. Now, let's review the Bayesian method once again.

% \section{ベイズ推定の基礎}
\section{Basics of Bayesian estimation}

% データ生成モデルが確率の言葉で記述される，というのはすでに述べたところです。推測統計というのがそもそも不良設定問題，すなわち少なすぎるヒントから未知なるものを当てるという状況に置かれた学問であり，この「わからないこと」を確率で表現するからです。この確率についての考え方として，ベイズの公式というのがあるのでした。ベイズの定理は次のように書き表されます。
I have already mentioned that the data generation model is described in terms of probability. Inferential statistics is a discipline placed in a situation of ill-posed problems, that is, guessing the unknown from too few hints, and expresses this "unknown" in terms of probability. There is a way of thinking about this probability, which is Bayes' theorem. Bayes' theorem is written as follows.

% $$ P(\theta|D) = \frac{P(D|\theta)P(\theta)}{P(D)} $$
As a translator, I can help with language translations. But the text you provided seems to be a statistical or mathematical formula, which doesn't need a language translation because it's universally used in its current format. If you want a specific part of that formula to be explained or if you want the context in which it is used to be translated, feel free to give me more details.

% ここで$P(X)$はXについての確率，という表現です。$P(A|B)$は条件付き確率というもので，Bが与えられた時のAの確率，という意味です。
Here, $P(X)$ represents the probability with respect to X. $P(A|B)$ is what is called conditional probability, meaning the probability of A given B.

% ベイズの定理の用語は次の通りです。まず右辺の分子に注目しましょう。$P(D|\theta)$とありますが，ここで$D$はデータを表しています。$\theta$はデータを生む確率のパラメータです。$P(D|\theta)$はパラメータ$\theta$のもとでのデータ$D$が得られる程度を表すもの，という意味になります。これは\keyterm{尤度}{ゆうど}{likelihood}と呼ばれるもので，パラメータの関数になっているのでした。たとえば正規分布からデータが生成されると考えるのであれば，手元のデータがパラメータ$\theta$から出てくる可能性がどれぐらいあるのか，を表現していると言えるわけです。この尤度関数は，データを生み出すメカニズムの表現そのものです。
The terminology of Bayes' theorem is as follows. First, let's focus on the numerator on the right side. It says $P(D| \theta)$, where $D$ represents the data, and $\theta$ represents the parameters of the probability that generates the data. $P(D|\theta)$ represents the degree to which data $D$ can be obtained under parameter $\theta$, which is called the likelihood. This likelihood is a function of the parameters. For example, if we consider that the data is generated from a normal distribution, it can be said to express how likely it is that the data in hand will come from parameter $\theta$. This likelihood function is the direct representation of the mechanism generating the data.

% 右辺の分子の第二項，$P(\theta)$はパラメータの確率です。正規分布からデータが生成される例で言えば，尤度はあるパラメータの値からデータが得られる程度を表現しているのですが，そのパラメータが出てくる確率はそもそもどの程度であるか，を考えているのです。これは別名\keyterm{事前分布}{じぜんぶんぷ}{prior distribution}とも呼ばれます。実際にデータが出てくる前の段階で，そもそもそのパラメータがどれほど出やすいかという確率を表現しており，これは今回のデータを取る前までの事前のデータや経験に基づいている，ということができます。
The second term of the numerator on the right side, $P(\theta)$, is the probability of the parameters. Speaking in the case where data is generated from a normal distribution, the likelihood is expressing the degree to which data can be obtained from a certain parameter value, but this is considering to some degree the probability that these parameters will occur in the first place. This is also known as the "prior distribution". The probability of how likely those parameters are to occur is expressed in reality before the data is actually produced, and it can be said that this is based on the pre-existing data or experience before the current data was collected.

% 右辺の分母，$P(D)$はデータ全体が得られる確率であり，\keyterm{周辺尤度}{しゅうへんゆうど}{marginal likelihood}とか\keyterm{エビデンス}{えびでんす}{evidence}と呼ばれます。\keyterm{正規化定数}{せいきかていすう}{normalized constant}ということもあります。これについては理解しにくいところもあるかもしれませんが，後に述べる理由によって，ひとまず深く考える必要はありません。
The denominator of the right-hand side, $P(D)$, is the probability that all the data is obtained and is called the \keyterm{marginal likelihood} or \keyterm{evidence}. It can also be referred to as the \keyterm{normalized constant}. You may find this difficult to understand, but for reasons we will explain later, you don't need to think too deep into it for now.

% 左辺はこの計算の結果得られる\keyterm{事後分布}{じごぶんぷ}{posterior distribution}を表しています。$P(\theta|D)$はデータ$D$で条件づけられたパラメータ$\theta$の確率です。我々は確率モデルを作るわけですが，そのパラメータがどうなっているかは事前にはわかりません。が，データを取ることで「データが与えられたのであれば，パラメータはこうだよ」というのがわかるわけです。あるいは事前にパラメータはこのあたりにあるのではないか，と経験上考えていたとしても，それがデータを取ることによって更新される(確信が強くなったり，違うかもしれないと思い直したり)，ということを意味しています。
The left-hand side represents the \keyterm{posterior distribution} obtained as a result of this calculation. $P(\theta|D)$ is the probability of the parameter $\theta$ conditioned on data $D$. We create a probability model, but we don't know what the parameters will be beforehand. However, by taking data, we can find out "if data is given, the parameter is like this". Or even if we thought based on experience that the parameter might be around here beforehand, it means that it will be updated by taking the data (confidence is strengthened, or reconsidered that it might be different).

% この式を言葉で表現し直すと，次のようになります。
If this formula is rephrased in words, it will be as follows.

% $$ \mbox{事後分布} = \frac{\mbox{尤度}\times \mbox{事前分布}}{\mbox{周辺尤度}} = \mbox{尤度} \times \frac{事前分布}{周辺尤度}$$
$$ \mbox{Posterior Distribution} = \frac{\mbox{Likelihood}\times \mbox{Prior Distribution}}{\mbox{Marginal Likelihood}} = \mbox{Likelihood} \times \frac{\mbox{Prior Distribution}}{\mbox{Marginal Likelihood}}$$

% \keytermL{尤度}{ゆうど}がメカニズムそのものだ，ということはすでに書きました。たとえば\keytermL{回帰分析}{かいきぶんせき}においても，尤度を計算できます。普通の回帰分析は，誤差が確率的に生じると考え，誤差以外のところは線形モデルで表現します。すなわち，従属変数$Y_i$に対して，独立変数$X_i$があったとすると，
I have already written that the likelihood is the mechanism itself. For example, you can calculate the likelihood even in regression analysis. The usual regression analysis considers that errors occur probabilistically, and other than errors are represented by a linear model. In other words, given an independent variable $X_i$ for the dependent variable $Y_i$,

\begin{equation*}
	\begin{array}{lll}
		Y_i & =    & \hat{Y}_i + e_i = \beta_0 + \beta_1 X_i + e_i \\
		e_i & \sim & N(0,\sigma)                                   \\
	\end{array}
\end{equation*}

% より，
Sorry, but your message doesn't seem to contain any Japanese text for me to translate. Could you please provide it?
% $$ Y_i \sim N(\beta_0 + \beta_1 X_i, \sigma) $$
As a language translator, I'm unable to translate the given text as it's LaTeX - a high-quality typesetting system; it includes features designed for the production of technical and scientific content, not Japanese language. However, I can definitely help explain what this mathematical expression means:

The statement "$ Y_i \sim N(\beta_0 + \beta_1 X_i, \sigma) $" represents a statistical linear regression model where each output Y_i is normally distributed with a mean that is a linear function of the input X_i, and a standard deviation sigma. In other words, for each input value X_i, we expect the output value Y_i to be normally distributed around the value of $\beta$0 + $\beta$1X_i.
% となるのでした。ここで$N(\mu,\sigma)$は平均$\mu$，標準偏差$\sigma$の正規分布を表し，データ$Y_i$が正規分布から生成されているというモデルになっています。回帰分析の場合は\keytermL{最尤法}{さいゆうほう}で未知数$\beta_0,\beta_1,\sigma$を求めたりしましたが\footnote{あるいは\keyterm{最小二乗法}{さいしょうにじょうほう}{Least Square method}で求めるのですが，これは幸い最尤法の結果と一致します。正規分布を使った線形回帰モデルの場合，データの記述統計的値と確率モデルによる推定値が一致するので，気づかないまま推測統計学に足を踏み入れてしまえるのでした}，その名の通りこのモデルが示す尤度を最大にする未知数の求め方を最尤法というのでした。
So it is. Here, $N(\mu,\sigma)$ represents a normal distribution with mean $\mu$ and standard deviation $\sigma$, and the data $Y_i$ follows a model that is generated from a normal distribution. In the case of regression analysis, we used the \keytermL{maximum likelihood method}{さいゆうほう} to find the unknowns $\beta_0,\beta_1,\sigma$\footnote{Or, we find it using the \keyterm{least squares method}{さいしょうにじょうほう}{Least Square method}, which fortunately coincides with the result of the maximum likelihood method. In the case of a linear regression model using a normal distribution, the descriptive statistical value of the data and the estimated value by the probabilistic model coincide, so we could get involved in inferential statistics without realizing it}, and the way to obtain the unknowns that maximizes the likelihood shown by this model was called the maximum likelihood method.

% ですが，尤度は確率を表すものではありません。確率とは非負の実数で，すべての確率を足し合わせる(あるいは積分する)と$1.0$にならなければなりませんが，尤度は足し合わせても1.0にならない数字なのです。ですから，「手元のデータがパラメータからでてくる\textbf{可能性}」という変な言い回しをしていました。「出てくる\textbf{確率}」とは言えないのですね。そこで，尤度を確率に置き換えたい，と考えたときに便利なのがベイズの定理なのでした。ベイズの定理は尤度に事前分布をかけ，周辺尤度で割るという操作によって，事後分布が得られる式，と読むこともできます。\underline{事後分布は\keytermL{確率分布}{かくりつぶんぷ}です}。尤度にある数字をかけて(事後の)確率分布にしていると考えるといいでしょう。
However, likelihood is not a probability. Probability is a non-negative real number, which all probabilities must add up to (or integrate) to $1.0$, but likelihood is a number that does not add up to $1.0$. Therefore, we have been using the strange phrase, "the \textbf{possibility} of the data at hand coming from the parameter". It cannot be said to be the "\textbf{probability}" of coming out. Hence, when we want to replace the likelihood with probability, Bayes' theorem is useful. Bayes' theorem can be interpreted as a formula that gives the posterior distribution by multiplying the likelihood by the prior distribution and dividing by the marginal likelihood. The \underline{posterior distribution is a \keytermL{probability distribution}}. It is a good idea to think that you are multiplying the likelihood by some number to convert it into a (posterior) probability distribution.
% ちなみに事前分布も確率分布ですが，\keytermL{周辺尤度}{しゅうへんゆうど}は確率ではありません。たとえば度数を総度数で割った相対頻度は確率と考えることができますが，それと同じように，分子をとある数字で割って，全体を1.0に整えるための定数なのです。\keytermL{正規化定数}{せいきかていすう}というのはそういう意味ですね。定数倍(正確には定数の逆数)をかけて大きさの調整をしているだけですので，ベイズの式はさらに次のように書き直すことができます。
By the way, the prior distribution is a probability distribution, but the marginal likelihood is not a probability. For instance, you can consider the relative frequency, which is the number of instances divided by the total number of instances, as a probability. In the same way, it's a constant to divide the numerator by a certain number and adjust the total to 1.0. That's what we mean by a normalization constant. It's just scaling by a constant (actually, the inverse of the constant), so you can rewrite Bayes' theorem even further like this.

% $$ \mbox{事後分布} \propto \mbox{尤度}\times \mbox{事前分布} $$
$$ \mbox{Posterior distribution} \propto \mbox{Likelihood} \times \mbox{Prior distribution} $$

% ここで$\propto$は比例する，という関係を表しています。最終的には事後分布の形が知りたいのですが，その形を決めているのは分子の尤度と事前分布だけだ，ということになります。
Here, $\propto$ represents a proportional relationship. Eventually, we want to know the shape of the posterior distribution, but what determines its shape is only the likelihood and the prior distribution in the numerator.

% ここで事前分布が\keyterm{一様分布}{いちようぶんぷ}{uniform distribution}であれば，事後分布の形状には影響せず，事後分布は尤度そのものの形を反映することになります。従来のデータ駆動型分析では，データに対して事前の仮定を一切置かないということでしたが，ベイズ流の分析でもそのことは同様に表現できるわけです。
If the prior distribution is a uniform distribution here, it will not affect the shape of the posterior distribution, and the posterior distribution will reflect the shape of the likelihood itself. In traditional data-driven analysis, there was no prior assumption about the data, but in Bayesian analysis, this can also be expressed in the same way.

% 以上がベイズの定理についての簡単な復習でした。しかしベイズの定理自体は，1740年代には明らかになっていたことであり，それがなぜ21世紀の今になって見直されてきたのでしょうか。それには色々な理由がありますが，その1つは最近まで「ベイズ流の分析は絵に描いた餅」だったからです。すなわち理屈ではこのような形になることは明らかだったのですが，実際に計算するのは難しいことが多々ありました。その問題を解決したのが\keyterm{マルコフ連鎖モンテカルロ法}{まるこふれんさもんてかるろほう}{Markov Chain Monte-Carlo method; MCMC}と呼ばれる方法です。
That was a brief review of Bayes' theorem. However, Bayes' theorem itself was established in the 1740s, so why has it been reevaluated in the 21st century? There are various reasons for this, one of which is that until recently, "Bayesian analysis was like a pie in the sky". That is, it was clear in theory that it would take this form, but in practice, there were many difficulties in calculating it. The method that solved this problem is called the Markov Chain Monte-Carlo method, or MCMC.

% \section{マルコフ連鎖モンテカルロ法}
\section{Markov Chain Monte Carlo Method}
\label{MCMC}
% ベイズの定理から，事前分布と尤度が分かれば，計算式を解いて事後分布の形を算出できます。しかし\keytermL{確率分布}{かくりつぶんぷ}の式が複雑になればなるほど，その計算はとても難しくなります。確率関数を複数組み合わせたり，求めるべきパラメータの数が増えて行ったりすると，とてもじゃないけど計算して答えを出すことができない，ということになります。そのせいもあって，ベイズ統計学は長らく実践には向かない手法だったのですが，\keyterm{マルコフ連鎖モンテカルロ法}{まるこふれんさもんてかるろほう}{Markov Chain Monte-Carlo method}，略して\keytermL{MCMC}{MCMC}が出てきてからその状況が一変しました。
From Bayes' theorem, if you know the prior distribution and likelihood, you can solve the equations to determine the shape of the posterior distribution. However, the more complex the equation of the probability distribution becomes, the more difficult the calculation becomes. If we combine multiple probability functions, or if the number of parameters to be determined increases, it becomes impossible to carry out the calculation and obtain the answer. For this reason, Bayesian statistics was considered unsuitable for practical use for a long time. However, the situation has drastically changed since the advent of the Markov Chain Monte Carlo method, abbreviated as MCMC.

% MCMC法は近似的な答えを探す1つの方法です。MCMCという名前は「マルコフ連鎖」と「モンテカルロ法」の2つのパートからできあがっています。マルコフ連鎖は，目標となる確率分布状態になるような推移規則を作る方法であり，モンテカルロ法は乱数を発生させるアルゴリズムです。この2つが合体することで，確率分布の形がわからなくてもサンプルが得られるような，乱数発生器を作ることができるようになったのです。
The MCMC method is one way to find an approximate answer. The name MCMC is derived from two parts: "Markov Chain" and "Monte Carlo Method". A Markov chain is a method of creating a transition rule to reach the desired probabilistic distribution state, and the Monte Carlo method is an algorithm for generating random numbers. By combining these two, even if the shape of the probability distribution is unknown, a random number generator can be created that allows samples to be obtained.

% ベイズ統計のゴールは事後分布を作ることです。事前分布を一様分布にして，尤度はベルヌーイ分布にする，といった簡単な場合ですととくに問題ないのですが，ある確率分布のパラメータがあって，さらにそれを生成する確率分布があるといった，階層的に入れ子になっているようなモデルが出てくると，「最終的な事後分布の形がわからない」という状況になります。正規分布とかポアソン分布といった，名前がついている確率分布であればその特徴がはっきりわかるのですが，それらを組み合わせて作られるものは名もなき合成関数になり，どんな形状をしているのか，まったく想像つかないことがあります。マルコフ連鎖を使うと，そのなもなき合成関数の形状をとにかく作り上げることができる，そんな数学的技術です。
The goal of Bayesian statistics is to create a posterior distribution. This is particularly not an issue in simple cases, like setting the prior distribution to a uniform distribution and the likelihood to a Bernoulli distribution. However, when we come across hierarchical models like a probability distribution with parameters, and another probability distribution generating these parameters, it becomes difficult to determine the final shape of the posterior distribution. With probability distributions that have a given name like the normal distribution or Poisson distribution, it's easy to understand their characteristics. But when different ones are combined, they result in anonymous composite functions, and it's entirely possible that we have no idea what their shape might be. Using a Markov Chain allows us to create the shape of this anonymous function - a mathematical technique that allows us to do just that.

% モンテカルロ法は乱数発生技術です。乱数というのは規則性のない数ですが，これを作るのはなかなか難しいものです。人間が0-9の数字を適当に書いていくと，知らず知らずのうちに均等でないパターンができてしまいます\footnote{嘘の538(ゴサンパチ)という標語があって，人間がふと思いつきで数字を作ろうとすると5，3，8が多くできてしまうと言われています。エビデンス(出典)がわからないので本当かどうかわかりませんが・・・。}。コンピュータに(たとえばRに)乱数があるじゃないか，と思われるかもしれませんが，コンピュータはあくまでも計算機でどこまでも合理的です。乱数によって動いているように見えますが，乱数に見えるような数字を生成するアルゴリズムがその中には埋め込まれていますから，正確には擬似乱数でしかありません。コンピュータの作る乱数は，もちろん人間が適当に思いつく乱数よりもより規則性が少ないですが，それでも基本的に
The Monte Carlo method is a random number generation technique. Although a random number is a number without regularity, it is quite difficult to produce. If human beings write numbers from 0 to 9 arbitrarily, an uneven pattern will unwittingly be created\footnote{There is a slogan called false 538 (go-san-road), and it is said that if people try to create numbers suddenly, 5, 3, 8 will be created more frequently. I don't know whether it's true or not because I don't know the evidence (source).}. You might think that there are random numbers in a computer (for example, R), but a computer is after all a calculating machine and it is rational in every way. It may look like it is moving with random numbers, but the algorithm that generates seemingly random numbers is embedded in it, so it is exactly a pseudo-random number. The random numbers made by a computer are, of course, less regular than the random numbers people think of arbitrarily, but still basically...
\input{x17_Bayesian_escape_02}
% というステップを反復($t=1,2,3\ldots$)することで生成するのです。こうした乱数列ができれば，それを正規分布の形だとか二項分布の形に当てはめて出力することは簡単なのです。このステップ・バイ・ステップの計算法としてある確率分布に従う乱数発生技術があり，これがマルコフ過程とひっついてできたのがMCMC法ということになります\footnote{余談ですが，スマホアプリなどで「ガチャを引く」というのも内部では乱数が生成されていて擬似的に「偶然あたりが出た」のを装っているに過ぎません。私はゲームなどをする時，擬似乱数に思いを馳せ「どこかの誰かに遊ばされている」と思うとやる気がなくなるので，ガチャを引くようなシーンは興醒めしてしまいます。やはりマルチプレイヤーゲームのように，人間が背後にいたほうがよっぽど意外な行動が見られますし，ゲームよりもリアルの世界の方が擬似的でない本物の偶然を楽しむことができておもしろいと思います。皆さんはどうお考えですか。}。
This is generated by repeating the step ($t=1,2,3\ldots$). Once such a random number sequence is created, it is easy to fit it into the shape of a normal distribution or a binomial distribution and output it. There is a step-by-step computing method that generates random numbers following a certain probability distribution, and this, combined with the Markov process, has led to the MCMC method\footnote{As an aside, when you "pull a gacha" on smartphone apps, etc., internally random numbers are being generated and it is just pretending to have "accidentally hit a jackpot". When I play games, I lose interest when I think about the pseudo-random numbers and feel like I'm being played by someone somewhere, so I am turned off by scenes where I pull a gacha. After all, in multiplayer games where humans are behind the scenes, you can see much more unexpected actions, and I think it's more interesting to enjoy the real, non-pseudo randomness of the real world than games. What do you all think?}.


% MCMC法は，ですから，どんな形の確率分布であっても乱数のサンプリングは生成できる，という技術なわけです。事後分布の関数の形，形状がわからなくてもそこからの乱数は作れるという技術であり，コンピュータ技術の性能が発展した今では大量の乱数を生成することも瞬時に行われます。
The MCMC method is a technique that can generate random sampling from any form of probability distribution. It's technology that can create randomness even from a posterior distribution function whose form or shape is unknown, and with the advancements in computer technology today, large amounts of random numbers can be generated in an instant.

% 乱数を用いるアプローチはいくつかの利点があります。たとえば\keytermL{確率分布}{かくりつぶんぷ}の平均である期待値など，分布の特徴を記述するための計算は積分を含むので解析的に解くのは知識も技術も必要です。しかしその確率分布から乱数を発生させると，その平均値を求めることでその近似値を得ることができます。確率分布の計算が記述統計の計算に置き換えられるのであり，また統計ソフトウェアにとって大量のデータの記述統計量を描くのは造作もないことなのです。乱数による近似値は，あくまでも近似値，近くて似ている値に過ぎませんが，精度を上げるためにはその乱数の数を増やしてやるだけで良いことになります。この点もいいですね。
Using a random number approach has several advantages. For example, calculations to describe the characteristics of a probability distribution, such as its mean or expected value, involve integration which requires both knowledge and technique to solve analytically. However, by generating random numbers from that probability distribution, we can obtain an approximation by calculating its mean. The calculation of the probability distribution is replaced by the calculation of descriptive statistics, and it's no big deal for statistical software to draw a large amount of descriptive statistic data. The approximated value obtained by random numbers is, after all, an approximation, a value that is merely close and similar. To increase the accuracy, you just need to increase the number of random numbers, which is another good point.

% また求めたいパラメータが複数ある場合，たとえば正規分布だと平均と標準偏差が未知数ですし，回帰分析では平均の中に切片と傾きといった未知数が入っているわけですが，このような場合の事後分布は同時確率空間ということになります。すなわち平均と標準偏差という2つの変数の場合でも，可視化するなら3次元空間が必要です(x軸に平均，y軸に標準偏差，z軸に確率密度)。こうした多次元空間において，あるパラメータだけについて期待値を計算したい，という場合はそれ以外のパラメータについては\keytermL{周辺化}{しゅうへんか}といって積分して全部の可能性をつぶしてしまう必要があるのですが，この計算も解析的にやるには実に大変なものです。しかし多次元の事後分布空間から生成された乱数があるなら，注目したい変数だけについての記述統計をすれば，他の変数を周辺化したこと同じになります。なんと便利なのでしょう。
In the case where there are multiple parameters you want to estimate, for example, in a normal distribution the mean and standard deviation are unknowns, and in regression analysis, the mean includes unknowns such as the intercept and slope, the posterior distribution in such a case becomes what is known as a joint probability space. In other words, even in the case of two variables such as the mean and standard deviation, a three-dimensional space is needed to visualize them (x-axis for the mean, y-axis for standard deviation, z-axis for probability density). In such a multidimensional space, if you only want to calculate the expected value for a certain parameter, you need to integrate all other parameters, a process known as "marginalization", to eliminate all possibilities. However, this calculation is really hard to do analytically. But if you have random numbers generated from a multi-dimensional posterior distribution space, by performing descriptive statistics for the variable of interest only, it becomes the same as having marginalized the other variables. How convenient it is.

% このように，乱数を使ったアプローチは計算上非常に有利な特徴を揃えています。乱数を大量に発生させられるような計算機のパワーは，最近のPCでしたら十分持っていますから，最近になってベイズ統計学やモデリングアプローチが生きてきたわけですね。さらにありがたいことに，事後分布から乱数を発生させるためのプログラミング言語ツールが登場したのも大きいでしょう。古典的にはBUGSというソフトウェアが，最近ではJAGSやStanといったのがそれで，\keyterm{確率的プログラミング言語}{かくりつてきぷろぐらみんぐげんご}{stochastic programming language}と呼ばれたりします。これらの言語では，尤度と事前分布をモデルとして表記し，それにデータを与えてやることで，事後分布からの乱数を生成します。いわば万能乱数発生器なのです。こうした環境が整ったことで，簡単に分析できるようになりました。
Thus, the approach using random numbers has very advantageous features computationally. The power of a computer to generate a large number of random numbers is ample even on recent PCs, and so it's only recently that Bayesian statistics and modeling approaches have come into their own. It's also significant that programming language tools for generating random numbers from the posterior distribution have emerged. Classical software like BUGS, and more recent ones like JAGS and Stan, are examples of this, often referred to as stochastic programming languages. In these languages, you can model the likelihood and prior distributions, and by feeding data into it, you can generate random numbers from the posterior distribution. You could call it a universal random number generator. With these environments in place, it has become easy to perform analyses.

% \section{乱数によるアプローチの例}
\section{Example of Approach by Random Numbers}
% それではMCMCを使った実践に入る前に，乱数を使って何ができるのか，Rで確かめてみましょう。
Before we move on to practical application using MCMC, let's verify what we can do with random numbers in R.

% \subsection{乱数による近似の例}
\subsection{Example of Approximation by Random Numbers}
% まずは馴染み深い\keyterm{正規分布}{せいきぶんぷ}{Normal distribution}の例からいきましょう。
Let's start with a familiar example of the "Normal Distribution".

% 正規分布に感するRの関数は\texttt{dnorm,pnorm, qnorm ,rnorm}などがあります。この\texttt{d,p,q,r}は他の確率分布を表す関数の前につける接頭語で，\texttt{d}は確率密度，\texttt{p}は累積確率,\texttt{q}はある累積確率になるときの確率点，そして\texttt{r}が乱数発生を意味しているのでした。
The functions in R related to the normal distribution are \texttt{dnorm, pnorm, qnorm, rnorm}, etc. These prefixes of \texttt{d, p, q, r} are used before functions representing other probability distributions. Where \texttt{d} represents probability density, \texttt{p} represents cumulative probability, \texttt{q} represents the probability point when it becomes a certain cumulative probability, and \texttt{r} stands for random number generation.

\begin{lstlisting}[language=c++,label={code:17_01},caption={乱数のコードの例}]
rm(list = ls())
set.seed(12345)
dnorm(0, mean = 0, sd = 1)
pnorm(0, mean = 0, sd = 1)
qnorm(0.6, mean = 0, sd = 1)
rnorm(10, mean = 0, sd = 1)
set.seed(12345)
rnorm(5, mean = 0, sd = 1)
\end{lstlisting}

% \paragraph{コード解説}
\paragraph{Code Explanation}
\begin{description}
  \item[1行目] 環境の初期化
  \item[2行目] 乱数発生開始点の設定
  \item[3行目] 標準正規分布の，確率点$x=0$における確率密度の計算
  \item[4行目] 標準正規分布の，確率点$x=0$までの累積確率
  \item[5行目] 標準正規分布の，累積確率60\%になるときの確率点
  \item[6行目] 標準正規分布から乱数を10個出力する
  \item[7行目] 乱数発生開始点の再設定
  \item[8行目] 標準正規分布から乱数を5個出力する
\end{description}


% このコードを使って確認しておいて欲しいところは，3つあります。まず\texttt{d,p,q,r}をつけたときの意味です。それぞれ正規分布のどういう数字を返しているのか，しっかり確認しておいてください。
There are three points I'd like you to verify using this code. First is the meanings when \texttt{d,p,q,r} are attached. Please make sure to understand which number is returned from the normal distribution for each one.
% 次に7行目にある乱数の発生です。これを実行すると，標準正規分布にし違う乱数が10個出力されます。乱数ですので，数字の並びに規則性はありません。バラバラの数字が出ていることを確認してください。
Next is the generation of random numbers in line 7. When executing this, ten random numbers differing from the standard normal distribution are output. Since it is random, there is no regularity in the sequence of numbers. Please confirm that a variety of numbers are produced.
% 最後に8行目，9行目の内容です。8行目は2行目と同じく，乱数の開始点を定めています。Rで出力される乱数は，規則性がない数字とは言え，計算によって算出している数字ですから規則性は拭いきれず，再現してしまいます。どこから計算を始めるか，という開始点は\keyterm{シード値}{しーどち}{seed value}と呼ばれ，ここに入力した数字が開始点になります。シード値の設定にとくに意味はなく，好きな数字にしていただいて結構です。ポイントは，任意の数字であっても同じ数字に設定すると，乱数はそこから計算して算出されますので，9行目で実行した5つの乱数は7行目の最初の5つと同じ数字が再現されているというところです。
This is related to the contents of lines 8 and 9. The eighth line, like the second one, determines the starting point of random numbers. The random numbers output in R, although they are numbers without regularity, cannot eliminate regularity because they are numbers calculated by calculation, and they are reproduced. The starting point, where to start the calculation, is called the "seed value", and the number entered here becomes the starting point. There's no particular significance to the setting of the seed value, you're free to choose any number. The point is, even if it is an arbitrary number, if it is set to the same number, the random number will be calculated from there. Therefore, the five random numbers executed on line 9 reproduce the same numbers as the first five on line 7.
% 規則性のない数字なので再現されては困ると思うかもしれませんが，科学計算のアプローチという意味では再現性が担保できることも重要なのです。
You may think it would be problematic to reproduce these irregular numbers, but it's also important that reproducibility can be guaranteed in the sense of a scientific calculation approach.

% では乱数を使う例をもう少し見てみましょう。次は\keyterm{ベルヌーイ分布}{べるぬーいぶんぷ}{Bernoulli distribution}について考えてみたいと思います。これはコイントスをして表が出るか，裏が出るかという0/1の離散的結果を生み出す分布です。
Let's take a closer look at examples using random numbers. Next, I would like to think about the Bernoulli distribution. This distribution produces discrete 0/1 results, such as whether a coin toss lands on heads or tails.
% 残念ながらRにはベルヌーイ分布の関数はなく\footnote{\texttt{extraDistr}パッケージを用いると，\texttt{bern}という関数名でベルヌーイ分布が使えます。}，二項分布の特殊例として使うことにします。\keyterm{二項分布}{にこうぶんぷ}{Binomial distribution}とは，N回コイントスをしてK回表が出る確率を表す分布ですが，これのコイントス回数が1回であればベルヌーイ分布と同じことになります。
Unfortunately, R does not have a function for the Bernoulli distribution \footnote{With the \texttt{extraDistr} package, you can use the Bernoulli distribution under the function name \texttt{bern}.}. We will use it as a special case of the binomial distribution. The \keyterm{binomial distribution}{Binomial distribution} is a distribution that represents the probability of getting heads K times by tossing a coin N times. If the number of coin tosses is once, it’s equivalent to the Bernoulli distribution.
\begin{lstlisting}[language=c++,label={code:17_02},caption={乱数のコードの例2}]
rbinom(10, size = 1, prob = 0.5)
rbinom(10, size = 1, prob = 0.3)
rbinom(10, size = 1, prob = 0.7)
# パッケージの利用
library(extraDistr)
rbern(10, prob = 0.5)
\end{lstlisting}

% \paragraph{コード解説}
\paragraph{Code Explanation}
\begin{description}
	\item[1行目] 二項分布の乱数を10個発生させる。確率は0.5で。
	\item[2行目] 二項分布の乱数を10個発生させる。確率は0.3で。
	\item[3行目] 二項分布の乱数を10個発生させる。確率は0.7で。
	\item[4-6行目] \texttt{extraDistr}パッケージによるベルヌーイ乱数例
\end{description}


% ここで確認して欲しいのは，コイントスを10回するとして，確率0.5の場合はほぼ半々，0.3の場合は半分以下，0.7の場合は半分以上表($1$)が出ている，ということです\footnote{正確には二項分布なので，1が出た回数が出力されているのです。試行数(\texttt{size})を1にしているので，1回中1回表が出た＝表($1$)が出た，と解釈しても，結果的に同じことであるというだけです。}。ここで確率を幾つにするかは，われわれが自由に設定できます。また乱数を幾つ発生させるかも自由です。これらの数字を色々変えて遊んでみましょう。
What we want to verify here is that when we toss a coin 10 times, with a probability of 0.5, we get almost half heads and half tails. If the probability is 0.3, we will get less than half heads, if it's 0.7, we will get more than half heads\footnote{To be precise, this is a binomial distribution, so the number of times we get heads (1) is the output. Since we set the number of trials (\texttt{size}) to 1, we can interpret it as getting heads (1) once out of once, which ultimately means the same thing.}. We can set the probability to whatever we want and freely generate as many random numbers as we want. Please feel free to play around with these numbers.

% たとえば乱数の数を100回，1000回，10000回と増やしたとします。その結果をすべて表示するのは大変ですが，その平均値を計算してみるとどうなるでしょうか。
For example, let's say we increase the number of random numbers to 100 times, 1000 times, 10000 times. It's quite labor-intensive to display all results, but what would happen if we calculate the average?
\begin{lstlisting}[language=c++,label={code:17_03},caption={乱数のコードの例3}]
N100 <- rbinom(n = 100, size = 1, prob = 0.5)
N1000 <- rbinom(n = 1000, size = 1, prob = 0.5)
N10000 <- rbinom(n = 10000, size = 1, prob = 0.5)
mean(N100)
mean(N1000)
mean(N10000)
\end{lstlisting}

% \paragraph{コード解説}
\paragraph{Code Explanation}
\begin{description}
	\item[1-3行目] 二項分布の乱数を100個，1000個，10000個発生させる。確率は0.5で。
	\item[4-6行目] それぞれの平均値を計算する
\end{description}

% これの実行結果は，私の環境では出力\ref{RS:out17_01}のようになりました。
The execution result of this, in my environment, appeared as shown in Output\ref{RS:out17_01}.

\begin{Rscreen}{乱数出力の結果}{out17_01}
  \begin{verbatim}
> mean(N100)
[1] 0.62
> mean(N1000)
[1] 0.52
> mean(N10000)
[1] 0.5003\end{verbatim}
\end{Rscreen}


% もともと確率を\texttt{prob = 0.5}と設定しましたから，理論的には0.5，つまり半々の確率で表($1$)が出たり裏($0$)が出たりするはずですが，最初の100回では62回なのでやや表が多めに出たようです。乱数ですので，そういうことはあります。ただ，この傾向も，1000回，100000回と繰り返していくと，ほぼ偶然による偏りはなく，半々の比率になっていくのが確認できると思います。このように，乱数が十分多ければ，確率分布の近似値として使うことができるというわけです。
Originally, the probability was set to \texttt{prob = 0.5}, so theoretically it should be 0.5, in other words, there should be a 50/50 chance of getting heads ($1$) or tails ($0$). However, for the first 100 times, it seems that heads came up 62 times, which is a little more than expected. However, this is a random number, so such things can happen. However, as this trend continues 1000, 10000 times and beyond, any bias due to chance should disappear and it should become a 50/50 ratio. It's safe to say that you can confirm this. Therefore, when there are many random numbers, it can be used as an approximate value of the probability distribution.

% ちなみに今回は平均をとりましたが，これは確率分布の\keyterm{期待値}{きたいち}{Expectation}を計算したこと同じになります。
By the way, this time we took the average, but this is the same as calculating the \keyterm{expected value}{きたいち}{Expectation} of the probability distribution.
% \subsection{乱数を用いた確率分布のプロット}
\subsection{Plotting probability distributions using random numbers}
% 乱数を使って近似できる，という例を可視化で見てみましょう。また正規分布を例にしてみたいと思います。
Let's visualize an example of how we can approximate using random numbers. Also, I would like to use the normal distribution as an example.
\begin{lstlisting}[language=c++,label={code:17_04},caption={確率分布の可視化1}]
x <- seq(-4, 4, 0.05)
plot(x, dnorm(x))
library(tidyverse)
ggplot(data.frame(x = c(-4, 4)), aes(x = x)) +
  stat_function(fun = dnorm)
\end{lstlisting}

% \paragraph{コード解説}
\paragraph{Code Explanation}
\begin{description}
  \item[1行目] -4から4までの範囲で，0.05刻みの数列を作る
  \item[2行目] 低水準プロット。各確率点の確率密度をプロットする
  \item[3行目] パッケージの読み込み
  \item[4-5行目] \texttt{ggplot}による美しい描画。\verb|stat_function|は関数の結果を図にするもの
\end{description}

% これは理論的な正規分布の形を図示するものです。美しいですね。でも同様のことが，乱数を使ってもできます。
This illustrates the shape of a theoretical normal distribution. It's beautiful, isn't it? But the same thing can be done using random numbers.

\begin{lstlisting}[language=c++,label={code:17_05},caption={確率分布の可視化2}]
N <- 1000
X <- rnorm(N, mean = 0, sd = 1) %>% as.data.frame()
ggplot(data = X, aes(x = .)) +
  geom_histogram()
\end{lstlisting}

% \paragraph{コード解説}
\paragraph{Code Explanation}
\begin{description}
  \item[1行目] 発生させる乱数の数
  \item[2行目] 乱数を発生させ，\texttt{data.frame}型に組み上げる
  \item[3-4行目] \texttt{ggplot}による美しい描画。\verb|geom_histogram|になっているところに注意
\end{description}

% このコードの結果描かれるカーブは，それほど美しいものではないかもしれません。しかし発生させる乱数の数が増えればどうでしょうか。どんどん形が似ていくことがわかると思います。
The curve drawn by this code may not be so beautiful. But what if the number of random numbers generated increases? I think you'll notice that the shapes are becoming more and more similar.
% ここでのポイントは，ヒストグラムを描いたらその稜線が\keytermL{確率分布}{かくりつぶんぷ}の形になること，乱数の数が増えれば十分な近似になることです。よく確認しておいてください。
The point here is that the shape of the histogram becomes the probability distribution after it is drawn, and if the number of random numbers increases, it will be a sufficient approximation. Please check it carefully.

% \subsection{乱数を用いた確率分布の要約}
\subsection{Summary of Probability Distribution Using Random Numbers}
% 最後に，確率分布の特徴を記述するコードの書き方について練習しましょう。
Finally, let's practice how to write code that describes the characteristics of a probability distribution.
% 乱数発生によるアプローチは，記述統計量で確率分布の特徴を記述できます。確率分布の期待値は，その算術計算の平均で良いのですが，中央値やパーセンタイル，分散や標準偏差も分布を記述する重要な指標です。
The approach by random number generation allows you to describe the characteristics of a probability distribution with descriptive statistics. The expected value of the probability distribution can be well-described with the average of its arithmetic calculation, but the median, percentile, variance, and standard deviation are also important indicators for describing the distribution.
% さらに確率分布における，確率密度が最も高くなる点は，「最も生じる確率が高い点」という意味で重要ですが，それを求めるためには特別な関数を書く必要があります。
In probability distribution, the point where the probability density is highest is important in the sense of being the point with the highest likelihood of occurring. However, to calculate this, a special function needs to be written.

% まずは\texttt{data.frame}型にした正規乱数を使って，記述統計量を算出する例を見てみましょう。
Let's start by looking at an example of calculating descriptive statistics using normally distributed random numbers converted into a \texttt{data.frame} type.

\begin{lstlisting}[language=c++,label={code:17_06},caption={確率分布の要約}]
N <- 10000
X <- rnorm(N, mean = 0, sd = 1) %>%
  as.data.frame() %>%
  dplyr::rename(val = 1)
X %>%
  summarise(
    Exp = mean(val),
    SD = sd(val),
    Median = median(val),
    U50 = quantile(val, prob = 0.50),
    U90 = quantile(val, prob = 0.90),
    L90 = quantile(val, prob = 0.10)
  )
\end{lstlisting}

% \paragraph{コード解説}
\paragraph{Code Explanation}
\begin{description}
	\item[1行目] 乱数を10万個作ることにします。
	\item[2行目] 標準正規分布の乱数発生
	\item[3行目] データフレーム型に整形
	\item[4行目] 変数名がつけられていないので，\texttt{val}という名前をつけることに。ここでの$1$は1列目を意味している。
	\item[6行目から] データフレーム\texttt{X}を使って計算
	\item[7行目] \texttt{summarise}関数は要約計算。変数名$=$計算式で表す。
	\item[8行目] 期待値を\texttt{mean}関数で算出
	\item[9行目] 確率分布の標準偏差を\texttt{sd}関数で算出\footnote{Rの\texttt{sd}関数は不偏分散の平方根を計算しており，厳密には標本標準偏差の計算で良いが，サンプルサイズが十分に大きいのでほとんど影響しません。}
	\item[10行目] 中央値を\texttt{median}関数で計算
	\item[11-13行目] パーセンタイル点を算出。50パーセンタイル点は中央値に同じ。
\end{description}


% また，確率密度が最も高くなる点を算出するための関数を自分で作ってみることにします。次のようにコードを入力してください。
Also, I will try to create a function myself to calculate the point where the probability density becomes the highest. Please enter the code as follows:
\begin{lstlisting}[language=c++,label={code:17_07},caption={確率分布の要約}]
map_estimation <- function(z) {
  density(z)$x[which.max(density(z)$y)]
}
# 試してみる
X %>%
  summarise(
    MAP = map_estimation(val)
  )

\end{lstlisting}

% \paragraph{コード解説}
\paragraph{Code Explanation}
\begin{description}
  \item[1行目] 関数の宣言。\verb|map_estimation|という関数名を作る。この関数は引数$z$を取る
  \item[2行目] \texttt{density}関数は与えられた数列の関数を考え，そのカーネル密度を計算します。\texttt{which.max}関数はその中でも最も密度の高い値を返すものです。
  \item[3行目] 関数の終わり
  \item[5-8行目] さきほどの例で試してみる。
\end{description}


% このようにして，確率密度関数の近似値が簡単な関数で求められます。
In this way, the approximate value of the probability density function can be easily obtained by using a simple function.

% 次回以降は，モデルに基づいて計算された事後分布から得られる乱数について，これらの関数を使ってその特徴を記述していくことになります。
From the next time onwards, we will use these functions to describe the characteristics of the random numbers obtained from the posterior distribution calculated based on the model.
% 関数や指標の意味をしっかりと理解しておいてください。
Please make sure to fully understand the meaning of functions and indicators.
% \section{課題}
\section{Task}
% 次の計算をするRコードを記述し，提出してください。提出はRスクリプトファイルでもRmdファイルでも構いませんが，どの課題に対するコードなのかがわかるよう，コメントや説明文を記入しておくこと。なお提出されたコード単体でバグがなく動くことが確認できないものは，未提出扱いになります。コードの書き方などわからないところがあれば，曜日別TAか小杉までメールで連絡し，指導を受けてください。
Please write and submit the following calculation in R code. Submissions can be in the form of an R script file or an Rmd file, but please include comments or explanatory text so that it's clear which assignment the code is for. If the submitted code can't be confirmed to run without bugs on its own, it will be treated as non-submitted. If you have any questions about how to write the code or anything else, please contact the day-specific TA or Kosugi by email and receive guidance.
\begin{enumerate}
	\item 平均50，標準偏差10の正規分布に従う乱数を10万点生成し，その平均値，中央値，標準偏差，15パーセンタイル，75パーセンタイル，2.5パーセンタイル，97.5パーセンタイルを算出してください。
	\item さきほど求めた正規乱数の記述統計量が，理論値の近似値になっていることを確認します。\texttt{dnorm,pnorm,qnorm}などを使って理論値を算出してください。
	\item 正規乱数のデータセットのうち，30以上60未満の値が含まれる割合を計算してください。またそれが理論値の近似値になっていることを確認するため，\texttt{dnorm,pnorm,qnorm}などを使って理論値を算出してください。
\end{enumerate}

\end{document}
